\scanpage{87}
\mspage{80}
% archive sheet number: 86
soit
\[
(I_\alpha)_{\alpha\in A},\quad I_\alpha\in\widetilde I,\quad\operatorname{card}A\leq\pi
\]
\deleted{\unsure{et considérons}} prouvons qu'il existe $I'\in\widetilde I$ qui majore les $I_\alpha$, i.e. qui contient $I'_0=\bigcup_{\alpha\in A}I_\alpha$. On aura $\operatorname{card}I'_0\leq\pi^2=\pi$ (car $\operatorname{card}A\leq\pi$ et $\operatorname{card}I_\alpha\leq\pi$ pour tt $\alpha$), et il suffit de voir ceci:

\uline{Lemme 1}: Soit \inserted{$\pi_0,\pi$ deux cardinaux, avec $\pi\geq\pi_0$,} $I$ un ens. ordonné grand devant $\pi$, et $I'_0$ une partie de $I$ de cardinal $\leq\pi_0$. Alors il existe une partie $I'$ de $I$, \deleted{contenant $I'_0$, et} satisfaisant les conditions
\[
\begin{aligned}
\text{a)}\quad&I'\supset I'_0\\
\text{b)}\quad&\operatorname{card}I'\leq\pi^{\ast}\,,\\
\text{c)}\quad&I'\text{ grand devant }\pi_0.
\end{aligned}
\]
% slanted; linked to the star beside the cardinal bound.
\margin{Pourvu qu'on suppose $\pi>\pi_0$, et de plus $\pi^{\pi_0}=\pi$ (p.ex. $\pi=2^c$ avec $c\geq\pi_0$)}
On se propose de construire $I'$ par récurrence transfinie, en définissant les $I'_\alpha\subset I^{\unsure{\deleted{\xi}|\xi|\deleted{\ast}}}$ ainsi:
\[
\begin{aligned}
I'_0&=I'_0\\
I'_{\alpha\deleted{\unsure{a|0}}}&=\bigcup_{\beta<\alpha}I'_\beta\quad\text{si }\alpha\text{ est un ordinal limite}
\end{aligned}
\]
$I'_{\alpha+1}$ déduit de $I'_\alpha$ ainsi: pour tte partie $J$ de $I'_\alpha$, telle que $\operatorname{card}J\leq\pi_0$, qui \unsure{n'est pas majorée} dans $I'_\alpha$, on choisit un \unsure{major.|majorant} \inserted{$i_J$} dans $I$ (existe, car $I$ grand devant $\pi$ donc dev. $\pi_0$), et $I'_{\alpha+1}$ \unsure{s'obtient} en adjoignant ces $i_J$ à $I'_\alpha$.
